\section{Spectral compactness and the matching lower bound}\label{sec:lower}

The upper bound uses only one trial direction. To prove that its coefficient
is sharp, we must identify the limit of every low-energy two-form and exclude
mass concentrated in the caps. We first make the metrics exactly conical,
apply conic spectral compactness at fixed auxiliary radii, and then remove
that modification by a multiplicative comparison of eigenvalues.

\subsection{Comparison without derivatives of the metrics}
\begin{lemma}\label{lem:metric-spectra}
Let $g,h$ be smooth metrics on a closed $n$-manifold and suppose
$(1-\delta)g\le h\le(1+\delta)g$, with $0<\delta<1/2$.
For every ordered positive Hodge eigenvalue in degree $p$,
\[
 K(\delta)^{-1}\lambda_{p,j}^+(g)
 \le\lambda_{p,j}^+(h)\le K(\delta)\lambda_{p,j}^+(g),
 \qquad K(\delta)=1+O_n(\delta).
\]
No derivatives of the metric comparison are required.
\end{lemma}
\begin{proof}
The squared $L^2$ norms of $p$-forms satisfy
\[
 a_p\|u\|_g^2\le\|u\|_h^2\le b_p\|u\|_g^2,
 \quad a_p=(1-\delta)^{n/2}(1+\delta)^{-p},\quad
 b_p=(1+\delta)^{n/2}(1-\delta)^{-p}.
\]
The distributional differential $d$ and its graph domain are unchanged
under equivalent norms. Set $Z^p=\ker d$ and give
$\operatorname{Dom}d/Z^p$ the quotient norm
$\|[u]\|_{g,q}^2=\inf_{z\in Z^p}\|u+z\|_g^2$.
The same constants compare these quotient norms. The coexact positive
$p$-eigenvalues are the squared singular values of $d$ on this quotient,
with min--max quotient $\|du\|_{p+1,g}^2/\|[u]\|_{g,q}^2$.
Its comparison factors are $a_{p+1}/b_p$ and $b_{p+1}/a_p$.
The positive exact $p$-spectrum equals the coexact $(p-1)$-spectrum.
The full positive spectrum is the ordered union of these two lists,
with multiplicity; the same common multiplicative bounds hold for their
ordered union. This proves the lemma and avoids comparing codifferentials
pointwise.
\end{proof}

\subsection{Fixed conic models}
On the cap end choose the radial variable $1/\rho_2$ from the product
corner in Section~\ref{sec:geometry}. In the notation of
\cite[\S5.3.2]{benabida}, its restriction is
$\sqrt{1+\widetilde r^2}$, where $\widetilde r$ is the original cap
radius. This changes that radius by a relative $O(\widetilde r^{-2})$
term, with differentiated bounds, and preserves the asymptotically
conical estimates after decreasing their positive decay rate if needed.
Fix a small central radius $R>0$ and a large cap radius $A>0$.
Let $g_0^R$ equal $g_0$ outside $r=2R$, equal $g_C$ inside $r=R$,
and interpolate between them. Let $g_{\mathrm{St}}^A$ agree with the
normalized Stenzel metric up to radius $A$, with $g_C$ beyond $2A$,
and interpolate on the intervening annulus. These are smooth positive
metrics; they need not be K\"ahler. Their relative errors are respectively
$O(R^\mu)$ and $O(A^{-a_1})$.

For $2A\epsilon<R$, glue the fixed central metric to the scaled metrics
$\epsilon^2g_{\mathrm{St}}^A$, using the actual link attachments. The
overlap is an exact cone. Fixed local normalization factors are absorbed
in the finitely many cap metrics. Denote the resulting smooth metric on
the same differentiable fibre by $\bar g_\epsilon^{R,A}$.

\begin{lemma}\label{lem:model-comparison}
For some $b>0$,
\begin{equation}\label{eq:model-comparison}
 \limsup_{\epsilon\downarrow0}\sup_{X_t}
 |g_\epsilon-\bar g_\epsilon^{R,A}|_{\bar g_\epsilon^{R,A}}
 \le C(R^\mu+A^{-b}).
\end{equation}
The family limit is taken with $R,A$ fixed.
\end{lemma}
\begin{proof}
On $r\ge R$, the actual metric converges to $g_0$, and the central
interpolation costs $O(R^\mu)$. On each fixed rescaled cap the actual
metric converges to Stenzel, and the cap interpolation costs
$O(A^{-a_1})$. In the intervening region $A\epsilon\le r\le R$,
Proposition~\ref{prop:metric-comparison} and
\eqref{eq:central-rate} give relative cone error
$O(r^\mu)+O((\epsilon/r)^{a_0})$, bounded by
$C(R^\mu+A^{-a_0})$. The corner coordinates agree at the two faces:
$r=\rho_1$ and the cap radius is $r/\epsilon=1/\rho_2$.
Thus these estimates compare the same conic frames. Taking
$b=\min(a_0,a_1)$ proves the assertion. The errors tending to zero
with $\epsilon$ may depend on $R,A$, which is why the limits are ordered.
\end{proof}

\subsection{Central harmonic forms and their domains}
\begin{lemma}\label{lem:central-kernels}
For $g_0^R$, the smooth $L^2$ closed and coclosed one-forms vanish.
The corresponding two-form space maps isomorphically to
$H^2(M_0;\RR)=\RR H\oplus\RR D$. Every such two-form lies in the
minimal graph domain of $d+d^*$.
\end{lemma}
\begin{proof}
On an exact conical end write a closed and coclosed two-form as
$u=a(r)+dr\wedge b(r)$. Link Hodge decomposition gives
\[
 a(r)=\sum_jc_j(r)\vartheta_j+d_Lv(r),\qquad d_L^*v=0,
\]
where $v$ is orthogonal to harmonic one-forms. The closedness equations
imply $c_j'=0$ and $b-v'$ closed. Coclosedness then gives
$b=v'+b_H$ with $b_H$ harmonic, and
\[
 r^{-2}\Delta_Lv-v''-3r^{-1}v'=0,
 \qquad b_H'+3r^{-1}b_H=0.
\]
The latter term is $r^{-3}B_H$; as a radial two-form it is not $L^2$
at zero unless zero. On a positive coexact one-form eigenspace of
link eigenvalue $\lambda$, the radial exponents of $v$ are
\[
 k_\pm=-1\pm\sqrt{1+\lambda}.
\]
The negative root is not $L^2$ as a two-form. Therefore
\begin{equation}\label{eq:cone-harmonic-expansion}
 u=\sum_jc_j\vartheta_j+dv_+,
 \quad |v_+|=O(r^{\kappa-1}),\quad |dv_+|=O(r^{\kappa-2})
\end{equation}
for some $\kappa>0$, with differentiated estimates. Orthogonal link
eigenfunction expansion and elliptic estimates on fixed-ratio annuli
justify these bounds; higher positive roots decay faster on a smaller
end. The harmonic link part admits the logarithmic cutoffs of
Proposition~\ref{prop:mass}, and the positive-root part admits ordinary
fixed-ratio cutoffs with vanishing graph-norm error. This proves the
minimal-domain assertion.

If $u$ is globally exact, choose a smooth primitive $w$. Its harmonic
link coefficients vanish. Since $H^1(T^{1,1})=0$, $w-v_+$ is locally
exact. Subtracting differentials of cutoff potentials makes $w=v_+$
near all ends. This primitive is in $L^2$ and in the minimal domain of
$d$: the squared cutoff error is $O(\rho^{2\kappa+2})$.
The adjoint identity then gives
$\|u\|^2=\langle d_{\min}w,u\rangle=\langle w,d_{\min}^*u\rangle=0$.
Thus the map to ordinary cohomology is injective.

For surjectivity, use the closed minimal de Rham complex. Its reduced
cohomology $\ker d_{\min}/\overline{\im d_{\min}}$ has harmonic
representatives by Hilbert-space orthogonal decomposition; no closed-range
assertion is needed for this reduced statement. The forms $\omega,\eta$
from Proposition~\ref{prop:mass} remain closed minimal-domain forms for
$g_0^R$, by norm equivalence. Their reduced classes are independent:
pair them with compactly supported closed four-forms dual to $H,D$
under $H^2(M_0)\times H_c^4(M_0)$. Those pairings are $L^2$-continuous
and vanish on the closure of minimal exact forms. The resulting reduced
harmonic representatives are smooth, closed and coclosed on the interior.
Injectivity into the two-dimensional ordinary cohomology now gives the
isomorphism.

For degree one, the cone calculation gives tangential harmonic link
one-forms and differentials of positive-root harmonic functions; the
radial solution $Cr^{-5}dr$ is not $L^2$. The positive scalar roots are
$-2+\sqrt{4+\lambda}>0$. Since $H^1(M_0)=0$, a harmonic one-form has
a global primitive that is bounded up to a constant on each end. Its
cutoff error has squared norm $O(\rho^4)$. The same adjoint argument
makes the form zero.
\end{proof}

Let $H_R,D_R$ be the harmonic representatives in the lemma and set
\begin{equation}\label{eq:model-mass}
 \alpha_R=D_R-\frac{\langle D_R,H_R\rangle}{\|H_R\|^2}H_R,
 \qquad M_R=\|\alpha_R\|^2.
\end{equation}
Then
\begin{equation}\label{eq:mass-limit}
 M_R\longrightarrow M\quad\text{as }R\downarrow0.
\end{equation}
Indeed $g_0^R$ and $g_0$ differ relatively by $O(R^\mu)$.
Their equivalent norms define the same closed minimal complex and
reduced cohomology. The harmonic norm of a class is the infimum over
its closed representatives modulo the closed exact subspace. Hence all
these norms differ by factors $1+O(R^\mu)$. Polarization on the
fixed two-dimensional space gives convergence of the Gram matrices,
and their Schur complements converge as well. The limiting harmonic
representatives are $\omega,\eta$ from Proposition~\ref{prop:mass}.
On the $i$th exact cone the harmonic link coefficient of $\alpha_R$
is the constant $b_i$, since $H$ has zero restriction there.

\subsection{Cap kernels and the critical degree}
We record the two facts needed to exclude cap mass in the spectral
compactness theorem. The complete cap with metric $g_{\mathrm{St}}^A$
is a scattering manifold after $x=1/r$. The scattering Hodge theorem
\cite[Theorem~1A]{hhm} identifies its $L^2$ harmonic forms below middle
degree with $H_c^j(T^*S^3)$, above middle degree with $H^j(T^*S^3)$,
and in middle degree with $\im(H_c^3\to H^3)$. All vanish. The first
and last ordinary groups vanish by retraction onto $S^3$ and duality;
the middle map is the Euler-class map of the cotangent bundle of $S^3$,
which is zero because $S^3$ is parallelizable. This conclusion does not
require the interpolated metric to be K\"ahler.

On a completed exact-cone cap, a \emph{critical extended solution} is a
smooth solution of $(d+d^*)u=0$ with expansion
$u=u_{\mathrm{crit}}+O(r^{-3-\kappa})$ at infinity for some $\kappa>0$,
with differentiated bounds. Here $u_{\mathrm{crit}}$ is homogeneous of
pointwise order $r^{-3}$; a nonzero such term has logarithmically divergent
squared $L^2$ norm. The critical coefficients of these global solutions,
modulo $L^2$ solutions, form the trace space $I_{1/2}$ in
\cite[equation~(8)]{anne-takahashi}. This is the space used in its
logarithmically normalized critical term. Throughout, physical degree
means ordinary differential-form degree.

\begin{lemma}\label{lem:critical-degree}
An actual critical extended solution of $d+d^*$ on a complete
six-dimensional exact-cone-ended cap has its logarithmic $L^2$ leading
term in physical degree three.
\end{lemma}
\begin{proof}
Such a solution has $|u|=O(r^{-3})$ and an expansion in powers from
the discrete tangential spectrum, with a positive gap after its critical
term. First, $du=d^*u=0$ separately, even if $u$ has mixed degree.
Choose a cutoff $\chi_T$ supported in $r<2T$, equal one below $T$,
with $|d\chi_T|\le C/T$. The shell mass is bounded, so
\[
 \|(d+d^*)(\chi_Tu)\|^2\le CT^{-2}.
\]
For compactly supported forms on a boundaryless manifold this is
$\|d(\chi_Tu)\|^2+\|d^*(\chi_Tu)\|^2$. On every fixed compact
set letting $T\to\infty$ proves the separate equations.

A homogeneous pure $p$-form of pointwise order $r^{-3}$ has the form
$u=r^{p-3}a+dr\wedge r^{p-4}b$. The leading closed/coclosed equations are
\[
 d_La=0,\quad (p-3)a=d_Lb,\quad d_L^*b=0,\quad
 d_L^*a+(p-3)b=0.
\]
If $p\ne3$, they imply $\Delta_Lb=-(p-3)^2b$, hence $a=b=0$.
For $p=3$ they state precisely that $a$ and $b$ are harmonic, yielding
$a+(dr/r)\wedge b$. Thus every critical leading term lies in degree
three. The argument applies to traces of actual complete-cap solutions,
not to arbitrary formal endpoint eigenvectors of a boundary operator.
The faster remainder is $L^2$ at infinity. After division by
$\sqrt{|\log\epsilon|}$, its contribution to the critical term
therefore tends to zero in $L^2$.
\end{proof}

\subsection{Counting the small eigenvalues}
\begin{proposition}\label{prop:model-count}
For each fixed $R,A$, the metrics $\bar g_\epsilon^{R,A}$ have a
uniform positive degree-one gap and exactly one positive degree-two
eigenvalue tending to zero. That eigenvalue is eventually simple and
coexact; the remaining positive degree-two spectrum has a uniform gap.
Its normalized real eigenforms converge, along subsequences and after
choice of sign, strongly on the central part and smoothly on fixed
compact subsets to $\alpha_R/\sqrt{M_R}$.
\end{proposition}
\begin{proof}
We apply the compactness and norm decomposition in
\cite[Theorems~A--B, Corollary~17, equations~(24)--(26), Lemma~18]{anne-takahashi}.
We retain the central extension selected by the cap trace space $W$
in \cite[equation~(7)]{anne-takahashi}. The additional link gap in
\cite[\S5.1, p.~1047]{anne-takahashi} identifies the minimal and
maximal Gauss--Bonnet domains and allows a cohomological description
of the full kernel. We do not use that identification. Instead,
Lemma~\ref{lem:central-kernels} places every relevant pure-degree
closed and coclosed limit in the minimal domain, which is contained
in the selected domain for every $W$.
Its geometric hypotheses are satisfied: the compact exterior is fixed
and exactly conical near its boundary, each compact cap has a conical
collar, and the caps are scaled by a common $\epsilon$, with fixed
normalizations. Several disjoint caps replace the boundary operator
and cap spaces by finite direct sums. All the local estimates in the
proof sum over those components; this introduces no independent
vanishing parameter.

For clarity, the consequence used here is the following norm
compactness statement. A normalized sequence of eigenforms with
vanishing eigenvalues has a subsequential central limit and possible
additional mass in complete-cap $L^2$ harmonic forms and critical
extended cap solutions. On the central region the sequence, up to an
$L^2$ error tending to zero, is the sum of a term converging strongly
to the central limit, a complete-cap tail, and a logarithmically
normalized critical term. The squared norm is the sum of the limiting
squared norms of those terms; orthogonality of finite collections is
preserved. The critical term is the radial extension of the trace of
an actual critical cap solution, with normalization
$|\log\epsilon|^{-1/2}$.

Apply this statement to pure $p$-form eigenvectors $u_j$, $p=1,2$,
with $\lambda_j\to0$. On the compact smooth manifolds,
\[
 \|du_j\|^2+\|d^*u_j\|^2=\lambda_j\longrightarrow0.
\]
The central limit is pure degree $p$, since it agrees with the local
limit of $u_j$ on every fixed compact subset. Testing the separate
energy bounds gives separate closedness and coclosedness. It is
therefore in the space of Lemma~\ref{lem:central-kernels}. This
argument makes no assertion that an arbitrary mixed-degree kernel of
a central Gauss--Bonnet extension is invariant under degree projection.

The complete-cap term vanishes by the $L^2$ calculation above. Write
the remaining three-term decomposition on the central part as
$u_j=\Phi_j+\widetilde\Phi_j+\overline\Phi_j+o_{L^2}(1)$.
Here $\Phi_j$ converges strongly to the central limit,
$\widetilde\Phi_j\to0$ by the cap-kernel vanishing, and
$\overline\Phi_j$ is the critical term. By
Lemma~\ref{lem:critical-degree} the latter has physical degree three,
up to an error tending to zero. Projecting this identity onto degree
three kills $u_j$ and its central limit; hence
$\|\overline\Phi_j\|\to0$. The norm identity now gives central norm
one, with no lost mass. This projection is applied to the actual
three-term decomposition, not commuted through its tangential spectral
projectors.

The central dimensions zero and two give the upper counts in degrees
one and two. Conversely approximate an orthonormal basis of the two
central harmonic two-forms by compactly supported pure two-forms in
the minimal graph domain. Their energies tend to zero and their Gram
matrix tends to the identity. For each fixed cutoff extend them by
zero across the shrinking caps. The model agrees with the central
metric on their supports when $\epsilon$ is small. Min--max, first
in $\epsilon$ and then in the cutoff radius, gives at least two
eigenvalues tending to zero, including the harmonic eigenvalue.
Since $b_1(X_t)=0$ and $b_2(X_t)=1$, precisely one positive two-form
eigenvalue tends to zero. Failure of either asserted complementary
gap would produce too many orthonormal zero-limit forms by compactness.
The same count implies eventual simplicity. The degree-one gap
excludes an exact small two-form, by compact Hodge decomposition.

It remains to identify its limit. Since $H$ restricts trivially to
the links, it has a compactly supported closed representative
$\gamma_H$ on $M_0$. Extend this by zero across the caps. Its model
harmonic representative has bounded norm by minimization, and a fixed
nonzero compactly supported dual pairing bounds its norm below.
The preceding compactness at eigenvalue zero gives convergence without
mass loss to $H_R$; compactly supported dual pairings identify its
class uniquely. The small eigenform is orthogonal to this harmonic
direction, so its central limit is $\pm\alpha_R/\sqrt{M_R}$.
Interior elliptic estimates upgrade convergence to smooth convergence
on fixed compact subsets. This proves the last assertion.
\end{proof}

\subsection{Sharp lower bound and proof of the main theorem}
Fix $R,A$ and a cone section $r_*=R/2$ on each end. Let
$B_{i,\epsilon}$ be the harmonic boundary coefficient of a normalized
small model eigenform on that section. Proposition~\ref{prop:model-count}
and \eqref{eq:cone-harmonic-expansion} imply, along every convergent
subsequence and up to common sign,
\[
 B_{i,\epsilon}\longrightarrow b_i/\sqrt{M_R}.
\]
Inside the section the metric is a scaled fixed exact-cone-ended cap.
Lemma~\ref{lem:capacity} therefore gives energy at least
\[
 \frac{cB_{i,\epsilon}^2}{\log(r_*/\epsilon)+C_{i,A}}.
\]
These regions are disjoint. Summing and using unit total norm yields
\begin{equation}\label{eq:model-liminf}
 \liminf_{\epsilon\downarrow0}
 \bar\lambda_2^+(\epsilon;R,A)\log(1/\epsilon)
 \ge \frac{c\sum_i b_i^2}{M_R}.
\end{equation}
The sign has disappeared on squaring. The argument applies to every
sequence by subsequence compactness, so it establishes the full lower
limit. The constants $C_{i,A}$ need no uniform bound in $A$, since the
family limit is taken first.

Lemmas~\ref{lem:metric-spectra} and~\ref{lem:model-comparison} transfer
\eqref{eq:model-liminf} to the actual metrics with factor tending to
one as $A\to\infty$ and $R\downarrow0$. By \eqref{eq:mass-limit},
\[
 \liminf_{\epsilon\downarrow0}\lambda_2^+(\epsilon)\log(1/\epsilon)
 \ge\frac{c\sum_i b_i^2}{M}.
\]
Proposition~\ref{prop:upper} is the matching upper limit. Substitution
of $c=2\pi/3$ and $\log(1/\epsilon)=|\log t|/3$ proves
\eqref{eq:main-law}, and \eqref{eq:epsilon-upper} proves
\eqref{eq:main-upper}.

For the gaps, fix a single admissible pair $R,A$ with small metric
comparison error. The model gaps transfer by a fixed positive factor.
Thus there is no need for a model gap uniform under the auxiliary
limits. Simplicity and coexactness follow as in
Proposition~\ref{prop:model-count}. On the actual K\"ahler manifold,
type and conjugation commute with the Laplacian. A nonzero $(2,0)$
component would supply two independent real eigenvectors, contradicting
simplicity. The eigenform is therefore real $(1,1)$. Its contraction
by $\Lambda_\omega$ is a scalar eigenfunction with the same positive
eigenvalue; if nonzero, its differential contradicts the degree-one
gap. The eigenform is primitive. This completes the proof of
Theorem~\ref{thm:main}.
